% GENERATED FILE. Edit canonical lessons, then run npm run generate:books.
% canonical-source-hash: fnv1a64:4b3668ddede13dc8
% canonical-lessons: SA-C82-lake, SA-C82-well, SA-C82-pond, SA-C82-island, SA-C82-shore

\chapter{Lake, Well and Shore}
\label{ch:sa-lake-well-shore}
\hlchaptermodality{82}

\begin{chapteropening}
\textbf{By the end of this chapter:} \emph{I can name a lake, a well, a pond, an island and a shore.}

Complete the last lesson of chapter 82: i can name a lake, a well, a pond, an island and a shore.
\end{chapteropening}

\section[saraḥ]{\sk{सरः} (saraḥ) — a lake}

\label{lesson:SA-C82-lake}



\begin{quote}
Before the new one: say the Sanskrit for a garden, then the Sanskrit for a forest.
\end{quote}

\subsection*{You'll want to know: \sk{सरः}}
\textbf{\sk{सरः}} (\emph{saraḥ}) — \textquotedblleft{}a lake\textquotedblright{}, still \textbf{\sk{जलम्}} where the lotus grows.

\begin{grammarlens}[title={What you've built}]
Still water.
\end{grammarlens}

\subsection*{Your turn}
\begin{itemize}
\raggedright
  \item \emph{Say it:} \emph{saraḥ}
  \item \emph{Say it:} \emph{saraḥ}, once more
  \item \emph{Say it:} say \emph{saraḥ}
  \item \emph{Recall:} say the Sanskrit for a garden, then the Sanskrit for a forest, then say \emph{saraḥ} again
\end{itemize}

\begin{tcolorbox}[breakable,colback=teal!4,colframe=teal!35!black,title={Before you move on}]
What does \sk{सरः} mean? (\textquotedblleft{}A lake\textquotedblright{}.) Say it once more.
\end{tcolorbox}
\section[kūpaḥ]{\sk{कूपः} (kūpaḥ) — a well}

\label{lesson:SA-C82-well}



\begin{quote}
Before the new one: say the Sanskrit for a forest, then the Sanskrit for a lake.
\end{quote}

\subsection*{You'll want to know: \sk{कूपः}}
\textbf{\sk{कूपः}} (\emph{kūpaḥ}) — \textquotedblleft{}a well\textquotedblright{}, water from under the ground.

\begin{grammarlens}[title={What you've built}]
Water from below.
\end{grammarlens}

\subsection*{Your turn}
\begin{itemize}
\raggedright
  \item \emph{Say it:} \emph{kūpaḥ}
  \item \emph{Say it:} \emph{kūpaḥ}, once more
  \item \emph{Say it:} say \emph{kūpaḥ}
  \item \emph{Recall:} say the Sanskrit for a forest, then the Sanskrit for a lake, then say \emph{kūpaḥ} again
\end{itemize}

\begin{tcolorbox}[breakable,colback=teal!4,colframe=teal!35!black,title={Before you move on}]
What does \sk{कूपः} mean? (\textquotedblleft{}A well\textquotedblright{}.) Say it once more.
\end{tcolorbox}
\section[taḍāgaḥ]{\sk{तडागः} (taḍāgaḥ) — a pond, a tank}

\label{lesson:SA-C82-pond}



\begin{quote}
Before the new one: say the Sanskrit for a lake, then the Sanskrit for a well.
\end{quote}

\subsection*{You'll want to know: \sk{तडागः}}
\textbf{\sk{तडागः}} (\emph{taḍāgaḥ}) — \textquotedblleft{}a pond\textquotedblright{}, a village tank with steps.

\begin{grammarlens}[title={What you've built}]
The village water.
\end{grammarlens}

\subsection*{Your turn}
\begin{itemize}
\raggedright
  \item \emph{Say it:} \emph{taḍāgaḥ}
  \item \emph{Say it:} \emph{taḍāgaḥ}, once more
  \item \emph{Say it:} say \emph{taḍāgaḥ}
  \item \emph{Recall:} say the Sanskrit for a lake, then the Sanskrit for a well, then say \emph{taḍāgaḥ} again
\end{itemize}

\begin{tcolorbox}[breakable,colback=teal!4,colframe=teal!35!black,title={Before you move on}]
What does \sk{तडागः} mean? (\textquotedblleft{}A pond, a tank\textquotedblright{}.) Say it once more.
\end{tcolorbox}
\section[dvīpaḥ]{\sk{द्वीपः} (dvīpaḥ) — an island}

\label{lesson:SA-C82-island}



\begin{quote}
Before the new one: say the Sanskrit for a well, then the Sanskrit for a pond.
\end{quote}

\subsection*{You'll want to know: \sk{द्वीपः}}
\textbf{\sk{द्वीपः}} (\emph{dvīpaḥ}) — \textquotedblleft{}an island\textquotedblright{}, land with water on both sides: \textbf{\sk{द्वि}}, \textquotedblleft{}two\textquotedblright{}.

\begin{grammarlens}[title={What you've built}]
Land in the water.
\end{grammarlens}

\subsection*{Your turn}
\begin{itemize}
\raggedright
  \item \emph{Say it:} \emph{dvīpaḥ}
  \item \emph{Say it:} \emph{dvīpaḥ}, once more
  \item \emph{Say it:} say \emph{dvīpaḥ}
  \item \emph{Recall:} say the Sanskrit for a well, then the Sanskrit for a pond, then say \emph{dvīpaḥ} again
\end{itemize}

\begin{tcolorbox}[breakable,colback=teal!4,colframe=teal!35!black,title={Before you move on}]
What does \sk{द्वीपः} mean? (\textquotedblleft{}An island\textquotedblright{}.) Say it once more.
\end{tcolorbox}
\section[tīram]{\sk{तीरम्} (tīram) — a shore, a bank}

\label{lesson:SA-C82-shore}



\begin{quote}
Before the new one: say the Sanskrit for a pond, then the Sanskrit for an island.
\end{quote}

\subsection*{You'll want to know: \sk{तीरम्}}
\textbf{\sk{तीरम्}} (\emph{tīram}) — \textquotedblleft{}a bank\textquotedblright{}, \textquotedblleft{}a shore\textquotedblright{}, the edge of a \textbf{\sk{नदी}} or the \textbf{\sk{समुद्रः}}.

\begin{grammarlens}[title={What you've built}]
The water's edge.
\end{grammarlens}

\subsection*{Your turn}
\begin{itemize}
\raggedright
  \item \emph{Say it:} \emph{tīram}
  \item \emph{Say it:} \emph{tīram}, once more
  \item \emph{Say it:} say \emph{nadītīram}
  \item \emph{Recall:} say the Sanskrit for a pond, then the Sanskrit for an island, then say \emph{tīram} again
\end{itemize}

\begin{tcolorbox}[breakable,colback=teal!4,colframe=teal!35!black,title={Before you move on}]
What does \sk{तीरम्} mean? (\textquotedblleft{}A shore, a bank\textquotedblright{}.) Say it once more.
\end{tcolorbox}
